\synctex=1
\documentclass[11pt,aspectratio=169]{beamer}

\usepackage[T1]{fontenc}
\usepackage[english,provide=*]{babel}
\usepackage{amsmath,amssymb,amsfonts}
\usepackage{mathtools}
\usepackage{bm}
\usepackage{tikz}
\usetikzlibrary{arrows.meta,positioning,shapes,calc}
\usepackage{pgfplots}
\usepackage{xcolor}
\usepackage{booktabs}
\pgfplotsset{compat=1.18}

\usefonttheme[onlymath]{serif}

\definecolor{darkblue}{RGB}{0,51,102}
\definecolor{lightblue}{RGB}{230,240,250}
\definecolor{darkgreen}{RGB}{0,102,51}
\definecolor{darkred}{RGB}{153,0,0}
\definecolor{orange}{RGB}{255,128,0}
\definecolor{violet}{RGB}{102,0,153}
\definecolor{lightgreen}{RGB}{220,245,220}
\definecolor{lightorange}{RGB}{255,240,220}
\definecolor{gray}{gray}{0.4}

% Git hash for the footer.
\usepackage{xstring}
\usepackage{catchfile}
\immediate\write18{git rev-parse HEAD > git.hash}
\CatchFileDef{\HEAD}{git.hash}{\endlinechar=-1}
\newcommand{\gitrevision}{\StrLeft{\HEAD}{7}}

\usepackage{ccicons}

\setbeamertemplate{footline}{
  \hfill\vspace*{1pt}\href{https://creativecommons.org/publicdomain/zero/1.0/deed.en}{\cczero}\enspace\href{https://git.ithaca.fr/stepan/continuo}{\gitrevision}\enspace--\enspace\insertframenumber/\inserttotalframenumber\enspace--\enspace\today\enspace}

\setbeamertemplate{navigation symbols}{}

\hypersetup{colorlinks=true}


\title{Nonlinear continuous-time New Keynesian models}
\author{Stéphane Adjemian}
% \date{May, 2026}
\institute{Université du Mans}

\AtBeginSection[]
{
  \begin{frame}
    \frametitle{Plan}
    \tableofcontents[currentsection]
  \end{frame}
}

\begin{document}

\begin{frame}
\titlepage
\end{frame}


\begin{frame}{Plan}
\tableofcontents
\end{frame}

\section{Introduction}

\begin{frame}{Introduction}

\begin{itemize}

\item This note provides algebraic derivations of nonlinear continuous-time New Keynesian models implemented in \href{https://continuo.adjemian.eu}{Continuo}.\newline

\item The three flavours differ in how habits enter the household's problem and the resulting marginal cost expressions.\newline

\item All models feature a zero lower bound on the nominal interest rate, Rotemberg price adjustment costs, and a temporary productivity boom as the exogenous shock.\newline

\item Code implementing these models is available in the \texttt{examples/nk-nonlinear/} directory of the \href{https://git.ithaca.fr/stepan/continuo}{Continuo} repository.

\end{itemize}

\end{frame}

%==============================================================================
\section{Common setup}
%==============================================================================

\begin{frame}{Common setup -- Overview}

\begin{itemize}

\item Three flavours of the fully nonlinear continuous-time New Keynesian
  model, all in the file \texttt{examples/nk-nonlinear/}:\newline

\begin{itemize}

\item \texttt{baseline.mod}: standard NK with no habits.\newline

\item \texttt{external\_habit.mod}: additive habits taken as given by
  the household (catching-up-with-the-Joneses).\newline

\item \texttt{internal\_habit.mod}: the same habits internalised, with
  a costate on the habit stock.

\end{itemize}

\item Household preferences, the firms' Rotemberg pricing problem, the
  monetary-policy rule with ZLB and the TFP shock are common to all
  three.\newline

\item The variants differ only in how the habit stock $X$ enters the
  household's problem.

\end{itemize}

\end{frame}

\begin{frame}{Common setup -- Households}

\begin{itemize}

\item Representative household maximises
  \[
  U \;=\; \int_0^\infty e^{-\rho t}\,\bigg[u\bigl(C(t),X(t)\bigr) -
  \frac{N(t)^{1+\eta}}{1+\eta}\bigg]\,dt,
  \]
  with discount rate $\rho>0$ and inverse Frisch elasticity~$\eta$.\newline

\item Felicity $u(C,X)$ depends on a habit stock $X$ (specialised
  below; $h=0$ removes it). Labour disutility is separable.\newline

\item Budget constraint, real one-period bonds $B$:
  \[
  \dot B(t) \;=\; \bigl(R(t)-\pi(t)\bigr)\,B(t) + w(t)\,N(t) - C(t),
  \]
  with $R$ the nominal policy rate, $\pi$ inflation, $w$ the real wage.

\end{itemize}

\end{frame}

\begin{frame}{Common setup -- Firms: demand and technology}

\begin{itemize}

\item A continuum of monopolistically competitive intermediate firms
  $i\in[0,1]$ supplies varieties of a single composite good. Each
  faces the constant-elasticity demand schedule
  \[
  y_i \;=\; \left(\frac{p_i}{P}\right)^{-\varepsilon} Y,
  \]
  with $P$ the aggregate price index, $Y$ aggregate output, and
  $\varepsilon>1$ the elasticity of substitution between varieties.\newline

\item Linear technology $y_i = A\,n_i$ (with $A=1$), the real
  marginal cost is the same across firms $MC = w/A$ .\newline

\item \textbf{Rotemberg (1982) price-adjustment cost.} Adjusting $p_i$
  at rate $\pi_i = \dot p_i / p_i$ costs $\phi/2 \cdot \pi_i^2$ per
  unit of output, i.e.\ a real flow cost of $\tfrac{\phi}{2}\,\pi_i^2\,Y$.\newline

\item The target for the real marginal cost is the inverse price markup,
  $MC^\star = (\varepsilon-1)/\varepsilon$. The NKPC will discipline
  $\dot\pi$ by the gap $MC - MC^\star$.

\end{itemize}

\end{frame}

\begin{frame}{Common setup -- Firms: optimisation}

\begin{itemize}

\item Working with the relative price $\xi_i = p_i/P$, its law of
  motion is
  \[
  \dot\xi_i \;=\; \xi_i\,(\pi_i - \pi),
  \qquad \pi = \dot P / P.
  \]

\item In equilibrium the firm is owned by the household, so it values
  future real profits at the household's stochastic discount factor
  \[
  \Lambda(t) \;\equiv\; u'(C(t)) \;=\; C(t)^{-\sigma},
  \]
  i.e.\ the marginal utility of an extra real unit of consumption.\newline

\item Firm $i$ chooses $\{\pi_i(t)\}$ to maximise
  \[
  \int_0^\infty e^{-\rho t}\,\Lambda(t)\,\Big[\,
    \underbrace{\xi_i^{1-\varepsilon}\,Y}_{\text{real revenue}}
    -
    \underbrace{MC\,\xi_i^{-\varepsilon}\,Y}_{\text{real cost}}
    -
    \underbrace{\tfrac{\phi}{2}\,\pi_i^2\,Y}_{\text{adj.\ cost}}
  \Big]\,dt
  \]
  subject to $\dot\xi_i = \xi_i\,(\pi_i - \pi)$, with $\xi_i(0)$
  given. The CES demand $y_i = \xi_i^{-\varepsilon} Y$ is substituted
  into the revenue and cost terms, so it does not appear as an
  explicit constraint.

\end{itemize}

\end{frame}

\begin{frame}{Common setup -- Firms: Hamiltonian}

\begin{itemize}

\item Casting the dynamic programme above as a Hamiltonian: state
  $\xi_i$, control $\pi_i$, costate $\nu$ on the relative-price
  constraint. The current-value Hamiltonian is
  \[
  \mathcal H \;=\; \Lambda\,\Big[\xi_i^{1-\varepsilon}\,Y
  - MC\,\xi_i^{-\varepsilon}\,Y
  - \tfrac{\phi}{2}\,\pi_i^2\,Y\Big]
  + \nu\,\xi_i\,(\pi_i - \pi).
  \]

\item Notice the structure: the bracket is the instantaneous real
  profit (weighted by the household's SDF $\Lambda$), and the
  $\nu\,\xi_i\,(\pi_i - \pi)$ term is the costate times the state's
  equation of motion --- the standard ``flow value + multiplier $\times$
  constraint'' form of an optimal-control Hamiltonian.\newline

\item The next slide derives the FOC in the control and the costate
  equation, then takes the symmetric-equilibrium limit.

\end{itemize}

\end{frame}

\begin{frame}{Common setup -- Firms: FOC and costate}

\begin{itemize}

\item FOC in the control $\pi_i$:
  \[
  \frac{\partial\mathcal H}{\partial\pi_i}
  = -\Lambda\,\phi\,\pi_i\,Y + \nu\,\xi_i = 0
  \quad\Longrightarrow\quad
  \nu\,\xi_i \;=\; \Lambda\,\phi\,\pi_i\,Y.
  \]

\item Costate equation
  $\dot\nu = \rho\,\nu - \partial\mathcal H / \partial\xi_i$, with
  \[
  \frac{\partial\mathcal H}{\partial\xi_i}
  = \Lambda\,\big[(1-\varepsilon)\,\xi_i^{-\varepsilon}\,Y
  + \varepsilon\,MC\,\xi_i^{-\varepsilon-1}\,Y\big]
  + \nu\,(\pi_i - \pi).
  \]

\item Impose symmetric equilibrium ($\xi_i = 1$, $\pi_i = \pi$, all
  firms identical). From the FOC, $\nu = \Lambda\,\phi\,\pi\,Y$. The
  last term of $\partial\mathcal H/\partial\xi_i$ vanishes, and the
  first two collapse to $\Lambda\,\varepsilon\,Y\,[MC - (\varepsilon-1)/\varepsilon]$:
  \[
  \dot\nu \;=\; \rho\,\nu - \Lambda\,\varepsilon\,Y\,
  \big[MC - \tfrac{\varepsilon-1}{\varepsilon}\big].
  \]

\item Differentiate $\nu = \Lambda\,\phi\,\pi\,Y$ exactly using the
  product rule:
  \[
  \dot\nu \;=\; \nu\,\Big[\frac{\dot\Lambda}{\Lambda}
  + \frac{\dot\pi}{\pi} + \frac{\dot Y}{Y}\Big].
  \]

\end{itemize}

\end{frame}

\begin{frame}{Common setup -- Firms: resource constraint and marginal cost}

\begin{itemize}

\item The Rotemberg adjustment cost is a \textbf{real resource loss}
  paid out of aggregate output. In symmetric equilibrium,
  \[
  Y \;=\; C + \tfrac{\phi}{2}\,\pi^2\,Y
  \quad\Longleftrightarrow\quad
  Y \;=\; \frac{C}{1 - (\phi/2)\,\pi^2}.
  \]

\item Differentiating $\log Y = \log C - \log(1 - (\phi/2)\pi^2)$
  yields the exact link
  \[
  \boxed{\frac{\dot Y}{Y} \;=\; \frac{\dot C}{C}
  \;+\; \frac{\phi\,\pi\,\dot\pi}{1 - (\phi/2)\,\pi^2}}
  \]
  used on a later slide to eliminate $\dot Y/Y$ from the costate
  identity.

\end{itemize}

\end{frame}

\begin{frame}{Common setup -- Firms: marginal cost across variants}

\begin{itemize}

\item Production $Y = A\,N$ gives $N = Y/A$. The labour FOC
  $w = N^\eta / u_C$ then writes the real marginal cost $MC = w/A$
  in terms of output and the model's particular marginal-utility-of-
  consumption expression $u_C$:
  \begin{align*}
  \text{baseline}\,(u_C = C^{-\sigma}): \qquad MC &= Y^\eta\,C^\sigma\,/\,A^{\eta+1}, \\
  \text{external habit}\,(u_C = (C-hX)^{-\sigma}): \qquad MC &= Y^\eta\,(C - h X)^\sigma\,/\,A^{\eta+1}, \\
  \text{internal habit}\,(u_C = \lambda_B - \lambda\mu): \qquad MC &= Y^\eta\,/\,(\lambda_B\,A^{\eta+1}).
  \end{align*}

\item At $A^\star = 1$ and $\pi=0$, $Y^\star = C^\star$ and these collapse to the
  canonical no-adjustment-cost expressions used in the steady-state
  formulas below.

\end{itemize}

\end{frame}

\begin{frame}{Common setup -- Firms: the exact NKPC, general $\sigma$}

\begin{itemize}

\item Combine $\dot\nu/\nu = \dot\Lambda/\Lambda + \dot\pi/\pi + \dot Y/Y$
  (slide before last) with $\dot\Lambda/\Lambda = -\sigma\,\dot C/C$ and
  the exact resource-constraint link
  $\dot Y/Y = \dot C/C + (\phi\pi\dot\pi)/(1-(\phi/2)\pi^2)$:
  \[
  \frac{\dot\nu}{\nu}
  \;=\; (1-\sigma)\,\frac{\dot Y}{Y}
  + \sigma\,\frac{\phi\,\pi\,\dot\pi}{1 - (\phi/2)\,\pi^2}
  + \frac{\dot\pi}{\pi}.
  \]

\item Equating with the costate result $\dot\nu/\nu = \rho - (\varepsilon/(\phi\pi))\,(MC - (\varepsilon-1)/\varepsilon)$ and multiplying by $\pi$:
  \[
  (1-\sigma)\,\frac{\dot Y}{Y}\,\pi
  + \frac{\sigma\,\phi\,\pi^2\,\dot\pi}{1 - (\phi/2)\,\pi^2}
  + \dot\pi
  \;=\; \rho\,\pi - \frac{\varepsilon}{\phi}\Big(MC - \tfrac{\varepsilon-1}{\varepsilon}\Big).
  \]


\end{itemize}

\end{frame}

\begin{frame}{Common setup -- Firms: the exact NKPC, log utility}

\begin{itemize}

\item Set $\sigma = 1$ (our calibration). The $(1-\sigma)\,\dot Y/Y\,\pi$
  term vanishes, and the previous slide's equation reduces to
  \[
  \dot\pi\,\Big[1 + \frac{\phi\,\pi^2}{1 - (\phi/2)\,\pi^2}\Big]
  \;=\; \rho\,\pi \;-\; \frac{\varepsilon}{\phi}\Big(MC - \tfrac{\varepsilon-1}{\varepsilon}\Big).
  \]

\item Using
  $1 + \phi\pi^2/(1-(\phi/2)\pi^2) = (1+(\phi/2)\pi^2)/(1-(\phi/2)\pi^2)$
  and dividing through:
  \[
  \boxed{\dot\pi \;=\; \frac{1 - (\phi/2)\,\pi^2}{1 + (\phi/2)\,\pi^2}\,
  \Bigl[\rho\,\pi - \tfrac{\varepsilon}{\phi}\bigl(MC - \tfrac{\varepsilon-1}{\varepsilon}\bigr)\Bigr]}
  \]

\item The correction factor is $1$ at $\pi = 0$.

\end{itemize}

\end{frame}

\begin{frame}{Common setup -- Monetary policy and the TFP shock}

\begin{itemize}

\item Taylor-type rule, clamped at the zero lower bound:
  \[
  \boxed{R(t) \;=\; \max\bigl(0,\ \rho + \phi_\pi\,\pi(t)\bigr)}
  \]
  Intercept $\rho$ so that the no-shock steady state has $R^\star=\rho$,
  $\pi^\star=0$. The $\max(0,\cdot)$ makes the model genuinely nonlinear.\newline

\item \textbf{TFP shock.} The only exogenous driver is total-factor
  productivity $A(t)$ in $Y = A\,N$.\newline

\item No-shock SS: $A^\star = 1$. The experiment is a temporary boom:
  $A$ rises to $1.12$ over $[0, 3)$ and returns to $1$. The boom lowers
  $MC$, drives deflation through the NKPC, and forces the Taylor rule
  onto the ZLB --- the classic ``supply-side deflationary trap''.

\end{itemize}

\end{frame}

\begin{frame}{Common setup -- Calibration}

\begin{center}
\small
\renewcommand{\arraystretch}{1.3}
\begin{tabular}{lcl}
\toprule
symbol & value & meaning \\
\midrule
$\sigma$    & $1$    & inverse IES (log utility) \\
$\eta$      & $1$    & inverse Frisch elasticity \\
$\varepsilon$ & $6$  & elasticity of substitution (20\% markup) \\
$\phi$      & $40$   & Rotemberg price-adjustment cost \\
$\rho$      & $0.02$ & rate of time preference (= steady-state real rate) \\
$\phi_\pi$  & $1.5$  & Taylor-rule inflation response \\
$\lambda$   & $0.5$  & habit adjustment speed (habit variants only) \\
$h$         & $0.7$  & habit weight (habit variants only) \\
\bottomrule
\end{tabular}
\end{center}

\end{frame}

%==============================================================================
\section{Baseline: no habit}
%==============================================================================

\begin{frame}{Baseline -- Euler equation}

\begin{itemize}

\item With $h=0$ the felicity collapses to
  $u(C) = C^{1-\sigma}/(1-\sigma)$, and the marginal utility of
  consumption is $u_C = C^{-\sigma}$.\newline

\item The continuous-time consumption Euler in a perfect-foresight
  setting reads $\dfrac{d\log u_C}{dt} = \rho - R + \pi$, so
  \[
  -\sigma\,\frac{\dot C}{C} = \rho - R + \pi,
  \]
  which rearranges to the standard form
  \[
  \boxed{\dot C \;=\; \frac{C}{\sigma}\,\bigl(R - \pi - \rho\bigr)}
  \]

\item For log utility ($\sigma=1$) this reduces to
  $\dot C/C = R - \pi - \rho$.

\end{itemize}

\end{frame}

\begin{frame}{Baseline -- Marginal cost}

\begin{itemize}

\item Labour FOC: $u_N = \lambda_B\,w$, with $u_N = -N^\eta$ (separable
  disutility) and $\lambda_B = u_C = C^{-\sigma}$ (no habit). Therefore
  \[
  w \;=\; \frac{N^\eta}{u_C} \;=\; N^\eta\,C^\sigma.
  \]

\item Goods-market clearing: $Y=C$ (production absorbs consumption,
  Rotemberg adjustment costs ignored at first order). Production gives
  $N = Y/A = C$ at $A=1$, so
  \[
  \boxed{MC \;=\; \frac{w}{A} \;=\; N^\eta\,C^\sigma \;=\;
  C^{\sigma+\eta}}
  \]

\item This closes the equation system together with the Euler and the
  NKPC.

\end{itemize}

\end{frame}

\begin{frame}{Baseline -- Full system}

\begin{center}
\small
\renewcommand{\arraystretch}{1.4}
\begin{tabular}{ll}
\toprule
role & equation \\
\midrule
state      & --- \\
jump       & $\dot C \;=\; (C/\sigma)\,(R-\pi-\rho)$ \\
jump       & $\dot\pi \;=\; \frac{1-(\phi/2)\pi^2}{1+(\phi/2)\pi^2}\,
              \bigl[\rho\,\pi - (\varepsilon/\phi)(MC - \tfrac{\varepsilon-1}{\varepsilon})\bigr]$ \\
algebraic  & $Y \;=\; C\,/\,(1 - (\phi/2)\,\pi^2)$ \\
algebraic  & $MC \;=\; Y^{\eta}\,C^{\sigma} / A^{\eta+1}$ \\
algebraic  & $R \;=\; \max(0,\ \rho + \phi_\pi\,\pi)$ \\
\bottomrule
\end{tabular}
\end{center}

\medskip

\textbf{Steady state} ($A^\star = 1$): $\pi^\star = 0$ gives $Y^\star = C^\star$, hence
\[
R^\star = \rho, \quad
MC^\star = \frac{\varepsilon-1}{\varepsilon}, \quad
C^\star = (MC^\star \cdot A^{*\,\eta+1})^{1/(\sigma+\eta)} = (MC^\star)^{1/(\sigma+\eta)}.
\]

\end{frame}

%==============================================================================
\section{External habit}
%==============================================================================

\begin{frame}{External habit -- Setup}

\begin{itemize}

\item Felicity is now
  \[
  u(C, X) \;=\; \frac{(C - h X)^{1-\sigma}}{1-\sigma},
  \qquad 0 \le h < 1,
  \]
  with a habit stock $X$ that evolves as an exponential moving average
  of consumption:
  \[
  \boxed{\dot X \;=\; \lambda\,(C - X), \qquad \lambda > 0}
  \]

\item \textbf{External}: the household treats $X$ as a given aggregate
  state — it does \emph{not} internalise that its own $C$ raises future
  $X$. There is no costate on $X$ in the household's problem.\newline

\item Marginal utility is
  \[
  u_C \;=\; (C - h X)^{-\sigma},
  \]
  exactly as in the no-habit problem but evaluated on the
  habit-adjusted consumption $C - h X$.

\end{itemize}

\end{frame}

\begin{frame}{External habit -- Euler equation}

\begin{itemize}

\item The Euler condition $\dfrac{d\log u_C}{dt} = \rho - R + \pi$ becomes
  \[
  -\sigma\,\frac{\dot C - h\,\dot X}{C - hX} \;=\; \rho - R + \pi.
  \]

\item Equivalently,
  \[
  \dot C - h\,\dot X \;=\; \frac{C - hX}{\sigma}\,(R - \pi - \rho),
  \]
  and substituting the habit law $\dot X = \lambda(C-X)$ yields the
  habit-augmented consumption Euler
  \[
  \boxed{\dot C \;=\; \frac{C - hX}{\sigma}\,(R - \pi - \rho)
  + h\,\lambda\,(C - X)}
  \]

\item When $h=0$ this reduces to the baseline Euler $\dot C =
  (C/\sigma)(R-\pi-\rho)$.

\end{itemize}

\end{frame}

\begin{frame}{External habit -- Marginal cost}

\begin{itemize}

\item The labour FOC is now
  \[
  w \;=\; \frac{N^\eta}{u_C} \;=\; N^\eta\,(C - hX)^\sigma.
  \]

\item Combining with $A=1$ and $N=C$ (goods-market clearing),
  \[
  \boxed{MC \;=\; \frac{w}{A} \;=\; C^\eta\,(C - hX)^\sigma}
  \]

\item In steady state $X^\star = C^\star$, so $C - hX = (1-h)C$ and
  $MC^\star = (1-h)^\sigma\,C^{*\,\sigma+\eta}$.

\end{itemize}

\end{frame}

\begin{frame}{External habit -- Full system}

\begin{center}
\small
\renewcommand{\arraystretch}{1.4}
\begin{tabular}{ll}
\toprule
role & equation \\
\midrule
state      & $\dot X \;=\; \lambda\,(C - X)$ \\
jump       & $\dot C \;=\; \tfrac{C - hX}{\sigma}\,(R - \pi - \rho)
              + h\,\lambda\,(C - X)$ \\
jump       & $\dot\pi \;=\; \frac{1-(\phi/2)\pi^2}{1+(\phi/2)\pi^2}\,
              \bigl[\rho\,\pi - (\varepsilon/\phi)(MC - \tfrac{\varepsilon-1}{\varepsilon})\bigr]$ \\
algebraic  & $Y \;=\; C\,/\,(1 - (\phi/2)\,\pi^2)$ \\
algebraic  & $MC \;=\; Y^\eta\,(C - hX)^\sigma / A^{\eta+1}$ \\
algebraic  & $R \;=\; \max(0,\ \rho + \phi_\pi\,\pi)$ \\
\bottomrule
\end{tabular}
\end{center}

\medskip

One state, two jumps, three algebraic relations — the same saddle-path
structure as the RBC, with the habit stock providing predetermined
persistence and $Y$ tracking $C$ up to the adjustment-cost wedge.

\end{frame}

\begin{frame}{External habit -- Steady state}

\begin{itemize}

\item $\dot X = 0 \Rightarrow X^\star = C^\star$, so $C - hX = (1-h)\,C^\star$ in
  steady state.\newline

\item $\dot\pi = 0$ and the Taylor rule give $\pi^\star = 0$, $R^\star = \rho$.\newline

\item Substituting into $MC^\star = (\varepsilon-1)/\varepsilon$:
  \[
  C^{*\,\eta}\,\bigl((1-h)\,C^\star\bigr)^\sigma \;=\;
  (1-h)^\sigma\,C^{*\,\sigma+\eta}
  \;=\; \frac{\varepsilon-1}{\varepsilon},
  \]
  hence
  \[
  \boxed{C^\star \;=\; \left(\frac{\varepsilon-1}{\varepsilon\,(1-h)^\sigma}\right)^{1/(\sigma+\eta)}}
  \]

\item When $h=0$ this collapses to the baseline steady state.

\end{itemize}

\end{frame}

%==============================================================================
\section{Internal habit}
%==============================================================================

\begin{frame}{Internal habit -- Hamiltonian}

\begin{itemize}

\item Same felicity $u(C,X) = (C - h X)^{1-\sigma}/(1-\sigma)$, same
  habit law $\dot X = \lambda(C - X)$, but the household now
  \textbf{internalises} the effect of $C$ on $X$.\newline

\item The current-value Hamiltonian carries a costate $\mu$ on the
  habit stock in addition to the wealth costate $\lambda_B$:
  \[
  \mathcal H \;=\; u(C, X) - \frac{N^{1+\eta}}{1+\eta}
  + \lambda_B\,\bigl[(R-\pi)\,B + w\,N - C\bigr]
  + \mu\,\lambda\,(C - X).
  \]

\item The first-order conditions and the two costate equations give
  the full system. $C$ is no longer chosen as a forward-looking variable
  — it is pinned algebraically at every instant by the consumption FOC.

\end{itemize}

\end{frame}

\begin{frame}{Internal habit -- First-order conditions}

\begin{itemize}

\item $\partial\mathcal H / \partial C = 0$:
  \[
  u_C - \lambda_B + \lambda\,\mu \;=\; 0
  \quad\Longleftrightarrow\quad
  \boxed{\lambda_B \;=\; (C - hX)^{-\sigma} + \lambda\,\mu}
  \]

  Compared with the no-habit case ($\lambda_B = u_C$), there is now an
  extra term $\lambda\,\mu$ capturing the marginal cost of raising the
  future habit stock.\newline

\item $\partial\mathcal H / \partial N = 0$: as before,
  \[
  -N^\eta + \lambda_B\,w \;=\; 0
  \quad\Longleftrightarrow\quad
  w \;=\; \frac{N^\eta}{\lambda_B}.
  \]
  With $Y=C=A\,N$ and $A=1$,
  \[
  \boxed{MC \;=\; \frac{w}{A} \;=\; \frac{C^\eta}{\lambda_B}}
  \]

\end{itemize}

\end{frame}

\begin{frame}{Internal habit -- Costate equations}

\begin{itemize}

\item \textbf{Wealth costate} $\lambda_B$:
  \[
  \boxed{\dot\lambda_B \;=\; \rho\,\lambda_B - \frac{\partial\mathcal H}{\partial B}
  \;=\; \lambda_B\,(\rho - R + \pi)}
  \]

\medskip
  
\item \textbf{Habit costate} $\mu$. With
  $u_X = -h\,(C - hX)^{-\sigma} = -h\,u_C$,
  \[
  \dot\mu \;=\; \rho\,\mu - \frac{\partial\mathcal H}{\partial X}
  \;=\; \rho\,\mu - \bigl(u_X - \mu\,\lambda\bigr)
  \;=\; (\rho + \lambda)\,\mu - u_X.
  \]
  Substituting $u_X$:
  \[
  \boxed{\dot\mu \;=\; (\rho + \lambda)\,\mu + h\,(C - hX)^{-\sigma}}
  \]

\end{itemize}

\end{frame}

\begin{frame}{Internal habit -- Why \texttt{C} is algebraic}

\begin{itemize}

\item The consumption FOC is an \emph{algebraic} relation between
  $\lambda_B$, $C$, $X$ and $\mu$:
  \[
  (C - hX)^{-\sigma} \;=\; \lambda_B - \lambda\,\mu.
  \]

\item Solving for $C$ (the unknown not constrained by any time
  derivative in this rearrangement):
  \[
  \boxed{C \;=\; h\,X + \bigl(\lambda_B - \lambda\,\mu\bigr)^{-1/\sigma}}
  \]

\item So in the BVP the predetermined state is $X$; the forward-looking
  jumps are $\pi$, $\lambda_B$ and $\mu$; the algebraic variables are
  $C$, $R$ and $MC$.\newline

\item This reorganisation is essential for the IR's classification rule:
  $\lambda_B$ has its own ODE and so cannot also be the LHS of an
  algebraic equation. Writing $C$ as the LHS of the FOC keeps each
  endogenous on exactly one side of one equation.

\end{itemize}

\end{frame}

\begin{frame}{Internal habit -- Full system}

\begin{center}
\small
\renewcommand{\arraystretch}{1.35}
\begin{tabular}{ll}
\toprule
role & equation \\
\midrule
state      & $\dot X \;=\; \lambda\,(C - X)$ \\
jump       & $\dot\lambda_B \;=\; \lambda_B\,(\rho - R + \pi)$ \\
jump       & $\dot\mu \;=\; (\rho + \lambda)\,\mu + h\,(C - hX)^{-\sigma}$ \\
jump       & $\dot\pi \;=\; \frac{1-(\phi/2)\pi^2}{1+(\phi/2)\pi^2}\,
              \bigl[\rho\,\pi - (\varepsilon/\phi)(MC - \tfrac{\varepsilon-1}{\varepsilon})\bigr]$ \\
algebraic  & $C \;=\; h\,X + (\lambda_B - \lambda\,\mu)^{-1/\sigma}$ \\
algebraic  & $Y \;=\; C\,/\,(1 - (\phi/2)\,\pi^2)$ \\
algebraic  & $MC \;=\; Y^\eta / (\lambda_B\,A^{\eta+1})$ \\
algebraic  & $R \;=\; \max(0,\ \rho + \phi_\pi\,\pi)$ \\
\bottomrule
\end{tabular}
\end{center}

\medskip

One state, three jumps, four algebraic equations. The terminal
boundary anchors $\pi$, $\lambda_B$ and $\mu$ at the steady state; the
initial state $X(0)$ is set by \texttt{initval(steady)}.

\end{frame}

\begin{frame}{Internal habit -- Steady state, part 1}

\begin{itemize}

\item $\dot\pi = 0$ and the Taylor rule give $\pi^\star = 0$ and
  $R^\star = \rho$.\newline

\item $\dot X = 0 \Rightarrow X^\star = C^\star$, so $C^\star - h X^\star = (1-h)\,C^\star$ in
  steady state.\newline

\item $\dot\mu = 0$:
  \[
  (\rho + \lambda)\,\mu^\star + h\,\bigl((1-h)\,C^\star\bigr)^{-\sigma} \;=\; 0
  \quad\Longleftrightarrow\quad
  \boxed{\mu^\star \;=\; -\,\frac{h\,\bigl((1-h)\,C^\star\bigr)^{-\sigma}}{\rho + \lambda} \;<\; 0}
  \]

  The shadow value of the habit stock is \emph{negative}: an extra unit
  of $X$ is a burden because it lowers future marginal utility.

\end{itemize}

\end{frame}

\begin{frame}{Internal habit -- Steady state, part 2}

\begin{itemize}

\item Consumption FOC at the steady state:
  \[
  \lambda_B^\star \;=\; (C^\star - h X^\star)^{-\sigma} + \lambda\,\mu^\star
  \;=\; \bigl((1-h)\,C^\star\bigr)^{-\sigma}\,\Bigl(1 - \frac{\lambda h}{\rho + \lambda}\Bigr),
  \]
  i.e.
  \[
  \boxed{\lambda_B^\star \;=\; \bigl((1-h)\,C^\star\bigr)^{-\sigma}\,
  \frac{\rho + \lambda - \lambda h}{\rho + \lambda}}
  \]

\item Substituting into $MC^\star = C^{*\,\eta}/\lambda_B^\star$ and equating to
  $(\varepsilon-1)/\varepsilon$:
  \[
  \boxed{C^\star \;=\; \left(\frac{\varepsilon-1}{\varepsilon\,(1-h)^\sigma}
  \cdot \frac{\rho + \lambda - \lambda h}{\rho + \lambda}\right)^{1/(\sigma+\eta)}}
  \]

\item At $h=0$ both correction factors vanish and $C^\star$ reduces to the
  baseline value $\bigl((\varepsilon-1)/\varepsilon\bigr)^{1/(\sigma+\eta)}$.

\end{itemize}

\end{frame}

%==============================================================================
\section{Numerical experiment}
%==============================================================================

\begin{frame}{Experiment -- The TFP boom}

\begin{itemize}

\item All three variants share the same shock specification:
  \[
  A(t) \;=\; 1 + 0.12\,\cdot\,\mathbf{1}_{[0,3)}(t),
  \]
  i.e.\ productivity rises 12\% over $[0, 3)$ and snaps back to
  $1$ thereafter. The path is known to agents at $t=0$ (single belief,
  one segment).\newline

\item Each model is solved on a uniform Crank--Nicolson grid with
  \texttt{simulate(T=25, N=600)}. The economy starts at its
  (model-specific) steady state through \texttt{initval(steady)} —
  vacuous for the baseline since it has no state, and pinning
  $X(0) = X^\star$ for the habit variants.\newline

\item The boom drives marginal cost down (firms can produce more
  cheaply); the NKPC delivers deflation; the Taylor rule wants to
  cut $R$ below zero but is clamped at the ZLB. \emph{What happens
  to consumption then depends sharply on preferences.}

\end{itemize}

\end{frame}

\begin{frame}{Experiment -- Headline results}

\begin{center}
\small
\renewcommand{\arraystretch}{1.3}
\begin{tabular}{lrrrr}
\toprule
scenario & $C$ on impact vs SS & $\pi$ trough & $R$ min & ZLB time \\
\midrule
baseline (no habit)        & $-4.7\%$ & $-6.60\%$ & $0$ & $9.8\%$ \\
external habit ($h=0.7$)   & $+7.0\%$ & $-5.00\%$ & $0$ & $8.0\%$ \\
internal habit ($h=0.7$)   & $+7.5\%$ & $-4.95\%$ & $0$ & $8.0\%$ \\
\bottomrule
\end{tabular}
\end{center}

\medskip

\begin{itemize}

\item \textbf{Baseline}: deflation and a binding ZLB produce a
  real-rate-driven recession ($C$ drops 4.7\% on impact) --- the
  classic supply-side deflationary trap.\newline

\item \textbf{Habit variants}: the wealth effect from the productivity
  boom dominates the trap channel. $C$ jumps \emph{up} 7\% on impact
  and decays smoothly back; deflation is milder; the ZLB still touches
  zero briefly.\newline

\item ``Trap mechanism'' and ``recession'' are not synonymous: the ZLB
  binds in all three cases, but the consumption response is sign-flipped
  by habit smoothing.

\end{itemize}

\end{frame}

\begin{frame}{Experiment -- Why the sign flip? Two competing channels}

\begin{itemize}

\item A TFP boom activates \textbf{two opposing forces} on consumption:
  \begin{itemize}
  \item \emph{Real-rate (substitution) channel}: deflation $+$ ZLB
    $\Rightarrow R - \pi > \rho$; the Euler
    $\dot c/c = (R-\pi-\rho)/\sigma$ pushes $c(0)$ \emph{down} to be
    consistent with $c(\infty)=c^\star$.\newline
  \item \emph{Wealth (income) channel}: higher $A$ raises lifetime
    resources, pushing $C(0)$ \emph{up}.
  \end{itemize}

\item \textbf{Baseline} (log utility, no habit). Income and substitution
  effects of the rate cancel exactly on the \emph{level} of $C$, so the
  Euler is the only pin-down. The trap delivers a high real rate, the
  substitution channel wins, and $C(0) < C^\star$.\newline

\item The next slide explains why habit \emph{reverses} this outcome
  by amplifying wealth and damping substitution.

\end{itemize}

\end{frame}

\begin{frame}{Experiment -- Why the sign flip? Habit amplifies wealth}

\begin{itemize}

\item \textbf{Habit amplifies the wealth channel.} Active consumption
  $c = C - hX$ has SS level $(1-h)C^\star = 0.3\,C^\star$; the same $\sigma$
  applied to a much smaller base means \emph{sharper curvature} in $C$.
  The household becomes more eager to grab transitory windfalls before
  habit catches up.\newline

\item \textbf{Habit dampens substitution.} The state $X$ moves only at
  rate $\lambda = 0.5$, so a one-shot consumption surge is \emph{not}
  immediately punished by a higher habit reference --- the utility cost
  is spread over years.\newline

\item \textbf{Internal habit makes the wealth channel explicit.} The
  consumption FOC $C = hX + (\lambda_B - \lambda\mu)^{-1/\sigma}$
  carries the forward-looking wealth costate $\lambda_B$, which
  integrates future productivity. The boom drops $\lambda_B$ on impact,
  and $C$ jumps up through the FOC --- the textbook permanent-income
  response, absent from the baseline because $C$ there has no level
  offset and is tied directly to the Euler.

\end{itemize}

\end{frame}

%==============================================================================
\section*{References}
%==============================================================================

\begin{frame}{References -- NK and monetary policy}

\begin{itemize}

\item Werning, I.\ (2011). Managing a Liquidity Trap: Monetary and
  Fiscal Policy. NBER Working Paper 17344.\newline

\item Rotemberg, J.J.\ (1982). Sticky Prices in the United States.
  \emph{Journal of Political Economy} 90(6), 1187--1211.\newline

\item Cochrane, J.H.\ (2017). The New-Keynesian Liquidity Trap.
  \emph{Journal of Monetary Economics} 92, 47--63.\newline

\item Gal\'i, J.\ (2015). \emph{Monetary Policy, Inflation, and the
  Business Cycle}. Princeton University Press.

\end{itemize}

\end{frame}

\begin{frame}{References -- Habit formation}

\begin{itemize}

\item Abel, A.B.\ (1990). Asset Prices under Habit Formation and
  Catching Up with the Joneses. \emph{American Economic Review Papers
  and Proceedings} 80(2), 38--42.\newline

\item Constantinides, G.M.\ (1990). Habit Formation: A Resolution of
  the Equity Premium Puzzle. \emph{Journal of Political Economy}
  98(3), 519--543.\newline

\item Christiano, L.J., Eichenbaum, M., and Evans, C.L.\ (2005).
  Nominal Rigidities and the Dynamic Effects of a Shock to Monetary
  Policy. \emph{Journal of Political Economy} 113(1), 1--45.

\end{itemize}

\end{frame}

\end{document}
